\section{O Conjunto de Dados: PublicHearingBR}
\label{sec:dados}

O PublicHearingBR ~\citep{fernandes2024publichearingbr} reúne 206 audiências públicas realizadas por
comissões da Câmara dos Deputados entre novembro de 2021 e maio de 2024. Cada amostra é uma tripla: a
transcrição oficial da sessão, a matéria publicada sobre ela pela Agência Câmara de Notícias e um resumo
estruturado com o assunto da audiência e as pessoas citadas na matéria, cada uma com nome, cargo e uma
lista de opiniões. O conjunto é aberto e está disponível no Hugging Face\footnote{\url{https://huggingface.co/datasets/unicamp-dl/PublicHearingBR}}.

As transcrições têm em média 18.102 palavras (mediana de 16.424 e máximo de 147.728), e as matérias,
627. Em média, a matéria tem 4,2\% do tamanho da transcrição e cita 5,2 participantes e 10,7 opiniões
\citep{fernandes2024publichearingbr}. No total, são 1.065 participantes e 2.203 opiniões.

O GPT-4o gerou, a partir da matéria e não da transcrição, uma versão preliminar do resumo estruturado
com 3.054 opiniões, e uma revisão manual contra a própria matéria corrigiu opiniões atribuídas à pessoa
errada e juntou opiniões divididas em duas, chegando às 2.203. Cada opinião registra, portanto, o que a
matéria atribui a uma pessoa, com a redação da matéria. O conjunto não indica em que ponto da
transcrição está a fala que originou cada opinião.

O dataset traz também um arquivo para inferência em linguagem natural (NLI), com 4.238 opiniões de
outra origem: elas foram geradas pelo ChatGPT a partir das transcrições, no experimento de sumarização
dos autores. Para cada opinião, um recuperador por \textit{embedding} selecionou os quatro trechos de
cinco sentenças mais próximos na fala da pessoa, e um consultor da Câmara anotou se a opinião pode ser
inferida desses trechos. Das 4.238, 504 (11,9\%) foram anotadas como não inferíveis. O rótulo se refere
aos quatro trechos, e não à transcrição inteira; por isso os autores tratam essa taxa como um limite
superior de alucinação.

\subsection{Pré-processamento e partição}
\label{sec:dados-particao}

Dividimos cada transcrição em turnos de fala a partir dos cabeçalhos que marcam cada troca de
falante, como ``\textit{A SRA. PRESIDENTE (Erika Kokay. PT - DF) -}'', o que resulta em 17.264 turnos
nas 206 audiências. O turno define o autor de
cada trecho da transcrição sem depender da matéria.

Como o arquivo não tem campo de data, usamos como data da audiência o carimbo de publicação da
matéria, presente nas 206. Nas 157
menções do tipo ``nesta quarta-feira (17)'' em que o dia da semana confere com o calendário, a
publicação ocorre no mesmo dia da audiência em 142 casos e no dia seguinte em 15. Com essas datas,
dividimos as audiências em ordem cronológica em treino, validação e teste, na proporção aproximada de
70/15/15 (Tabela~\ref{tab:dados}). Os cortes só foram feitos entre datas separadas por dois dias ou
mais, para que o erro de um dia não mude uma audiência de partição. A divisão temporal não separa
pessoas: 24 dos 879 participantes distintos citados nas matérias aparecem tanto no treino quanto no
teste. Na UDV (Seção~\ref{sec:metodo-udv}), o corte de similaridade é calibrado no treino, e a amostra de
validação humana é sorteada no teste (Seção~\ref{sec:setup-udv}).

Cada perfil por ator usado na avaliação (Seção~\ref{sec:metodo-perfis}) é escrito apenas com as
falas do ator nas audiências de treino, e as opiniões das audiências de teste servem para avaliá-lo.
As audiências de validação não entram nos perfis e, na simulação, servem apenas para escolher o
número de trechos recuperados (Seção~\ref{sec:setup-simulacao}). Aqui o ator é identificado pelo cabeçalho da
transcrição, e não pela citação na matéria, então conta também quem falou numa audiência sem ser
citado na matéria dela. Dos 301 atores que falam em duas ou mais audiências, 264 falam em ao menos uma
de treino, com 4.598 turnos em 139 das 144 audiências de treino. No teste, 106 das 359 opiniões se
ligam, pelo turno de fala da evidência da UDV, a 48 desses 264 atores, em 27 das 30 audiências.

\input{tables/dataset-statistics}
